
% ----------------------------------------------------------------------
%                   LATEX TEMPLATE FOR PhD THESIS
% ----------------------------------------------------------------------

% based on Harish Bhanderi's PhD/MPhil template, then Uni Cambridge
% http://www-h.eng.cam.ac.uk/help/tpl/textprocessing/ThesisStyle/
% corrected and extended in 2007 by Jakob Suckale, then MPI-CBG PhD programme
% and made available through OpenWetWare.org - the free biology wiki


%: Style file for Latex
% Most style definitions are in the external file PhDthesisPSnPDF.
% In this template package, it can be found in ./Latex/Classes/
\documentclass[twoside,11pt]{Latex/Classes/PhDthesisPSnPDF}


%: Macro file for Latex
% Macros help you summarise frequently repeated Latex commands.
% Here, they are placed in an external file /Latex/Macros/MacroFile1.tex
% An macro that you may use frequently is the figuremacro (see introduction.tex)
% This file contains macros that can be called up from connected TeX files
% It helps to summarise repeated code, e.g. figure insertion (see below).

% insert a centered figure with caption and description
% parameters 1:filename, 2:title, 3:description and label
\newcommand{\figuremacro}[3]{
	\begin{figure}[htbp]
		\centering
		\includegraphics[width=1\textwidth]{#1}
		\caption[#2]{\textbf{#2} - #3}
		\label{#1}
	\end{figure}
}

% insert a centered figure with caption and description AND WIDTH
% parameters 1:filename, 2:title, 3:description and label, 4: textwidth
% textwidth 1 means as text, 0.5 means half the width of the text
\newcommand{\figuremacroW}[4]{
	\begin{figure}[htbp]
		\centering
		\includegraphics[width=#4\textwidth]{#1}
		\caption[#2]{\textbf{#2} - #3}
		\label{#1}
	\end{figure}
}

% inserts a figure with wrapped around text; only suitable for NARROW figs
% o is for outside on a double paged document; others: l, r, i(inside)
% text and figure will each be half of the document width
% note: long captions often crash with adjacent content; take care
% in general: above 2 macro produce more reliable layout
\newcommand{\figuremacroN}[3]{
	\begin{wrapfigure}{o}{0.5\textwidth}
		\centering
		\includegraphics[width=0.48\textwidth]{#1}
		\caption[#2]{{\small\textbf{#2} - #3}}
		\label{#1}
	\end{wrapfigure}
}

% predefined commands by Harish
\newcommand{\PdfPsText}[2]{
  \ifpdf
     #1
  \else
     #2
  \fi
}

\newcommand{\IncludeGraphicsH}[3]{
  \PdfPsText{\includegraphics[height=#2]{#1}}{\includegraphics[bb = #3, height=#2]{#1}}
}

\newcommand{\IncludeGraphicsW}[3]{
  \PdfPsText{\includegraphics[width=#2]{#1}}{\includegraphics[bb = #3, width=#2]{#1}}
}

\newcommand{\InsertFig}[3]{
  \begin{figure}[!htbp]
    \begin{center}
      \leavevmode
      #1
      \caption{#2}
      \label{#3}
    \end{center}
  \end{figure}
}


%%% Local Variables: 
%%% mode: latex
%%% TeX-master: "~/Documents/LaTeX/CUEDThesisPSnPDF/thesis"
%%% End: 




%: ----------------------------------------------------------------------
%:                  TITLE PAGE: name, degree,..
% ----------------------------------------------------------------------
% below is to generate the title page with crest and author name

%if output to PDF then put the following in PDF header
\ifpdf  
    \pdfinfo { /Title  (PhD and MPhil Thesis Classes)
               /Creator (TeX)
               /Producer (pdfTeX)
               /Author (YourName your@email.net)
               /CreationDate (D:YYYYMMDDhhmmss)  %format D:YYYYMMDDhhmmss
               /ModDate (D:YYYYMMDDhhmm)
               /Subject (xyz)
               /Keywords (add, your, keywords, here) }
    \pdfcatalog { /PageMode (/UseOutlines)
                  /OpenAction (fitbh)  }
\fi


\title{Divergence of an introduced population of the
  Swimbladder-nematode \textit{Anguillicola crassus} - a
  transcriptomic perspective}



% ----------------------------------------------------------------------
% The section below defines www links/email for author and institutions
% They will appear on the title page of the PDF and can be clicked
\ifpdf
  \author{\href{mailto:emanuelheitlinger@gmail.com}{Emanuel Heitlinger}}
  

\cityofbirth{Schw\"abisch Gm\"und} % uncomment this if your university requires this
%  % If city of birth is required, also uncomment 2 sections in PhDthesisPSnPDF
%  % Just search for the "city" and you'll find them.
\collegeordept{\href{http://www.chem-bio.kit.edu/}{Fakult\"at für Chemie und Biowissenschaften}}
\university{\href{http://www.kit.edu/index.php}{Karlsruher Institut
    f\"ur Technologie (KIT) - Universit\"atsbereich}}

  % The crest is a graphics file of the logo of your research institution.
  % Place it in ./0_frontmatter/figures and specify the width
\crest{\includegraphics[width=4cm]{kit_logo_de.jpg}}
  
% If you are not creating a PDF then use the following. The default is PDF.
\else
  \author{Emanuel Heitlinger}
  \cityofbirth{born in Schw\"abisch Gm\"und}
  \collegeordept{Fakult\"at f\"ur Chemie und Biowissenschaften}
  \university{Karlsruher Institut f\"ur Technologie (KIT) - Universit\"atsbereich}
  \crest{\includegraphics[width=4cm]{kit_logo_de.jpg}}
\fi

%\renewcommand{\submittedtext}{change the default text here if needed}
\degree{Dr. rer. nat.}
\degreedate{2012 Januar}


% ----------------------------------------------------------------------
       
% turn of those nasty overfull and underfull hboxes
\hbadness=10000
\hfuzz=50pt


%: --------------------------------------------------------------
%:                  FRONT MATTER: dedications, abstract,..
% --------------------------------------------------------------

\begin{document}

%\language{english}

% sets line spacing
\renewcommand\baselinestretch{1.2}
\baselineskip=18pt plus1pt


%: ----------------------- generate cover page ------------------------

\maketitle  % command to print the title page with above variables


%: ----------------------- cover page back side ------------------------
% Your research institution may require reviewer names, etc.
% This cover back side is required by Dresden Med Fac; uncomment if needed.
%% Is empty the reviewers are on frontpage
\newpage

\newpage


%: ----------------------- abstract ------------------------

% Your institution may have specific regulations if you need an abstract and where it is to be placed in the document. The default here is just after title.


% Thesis Abstract -----------------------------------------------------


%\begin{abstractslong}    %uncommenting this line, gives a different abstract heading
\begin{abstractslong}        %this creates the heading for the abstract page

  The ability to expand into new environments and niches, despite
  being highly adapted to a habitual environment, is a fascinating
  feat of organisms. 30 years ago \textit{Anguillicola crassus} was
  introduced from Asia, where it parasitises \textit{Angilla
    japonica}, to Europe and spread here in the new host species
  \textit{Anguilla anguilla}. Whether and how much phenotypic
  plasticity or rapid adaptation to differential selection are
  contributing to its success in invading new host-populations is a
  question of substantial evolutionary interest.

  Gene regulatory networks, as an important link between genotype
  and phenotype, are thought to play a central role both in the
  response to stress (e.g. from as yet unexperienced environmental
  stressors) and in local adaptation.

  In the present project, differential gene-expression in
  \textit{A. crassus} populations was assessed using next generation
  sequencing on the 454 and Illumina platforms and genetic components
  of differences were isolated in cross-inoculation experiments with
  both Asian and European host-species and parasite populations.

  Several proteases were shown to be under positive selection on the
  sequence level, highlighting this group of enzymes as possible
  targets of an immune-attack on \textit{A. crassus}. On the
  gene-expression level, the extent of heritable change was large in
  comparison to the effect of modification in different
  host-environments. Mitochondrially encoded subunits of the
  respiratory chain and other genes connected to aerobic respiration
  showed divergent expression patterns in European vs. Asian parasite
  populations; cuticle collagen genes showed ``adapted'' patterns of
  expression in present day sympatric host-parasites pairs.
  
  These results identified gene-expression phenotypes, confirming the
  divergence of European \textit{A. crassus} populations. Such
  phenotypes will be more accessible to popluation-genetic analysis
  invetigating selection than complex life history traits.

\end{abstractslong}
%\end{abstractlongs}


% ---------------------------------------------------------------------- 

%%% Local Variables: ***
%%% mode:latex ***
%%% TeX-master: "../thesis.tex"  ***
%%% tex-main-file: "../thesis.tex" ***
%%% End: ***
     

% The original template provides and abstractseparate environment, if your institution requires them to be separate. I think it's easier to print the abstract from the complete thesis by restricting printing to the relevant page.
% \begin{abstractseparate}
%   
% Thesis Abstract -----------------------------------------------------


%\begin{abstractslong}    %uncommenting this line, gives a different abstract heading
\begin{abstractslong}        %this creates the heading for the abstract page

  The ability to expand into new environments and niches, despite
  being highly adapted to a habitual environment, is a fascinating
  feat of organisms. 30 years ago \textit{Anguillicola crassus} was
  introduced from Asia, where it parasitises \textit{Angilla
    japonica}, to Europe and spread here in the new host species
  \textit{Anguilla anguilla}. Whether and how much phenotypic
  plasticity or rapid adaptation to differential selection are
  contributing to its success in invading new host-populations is a
  question of substantial evolutionary interest.

  Gene regulatory networks, as an important link between genotype
  and phenotype, are thought to play a central role both in the
  response to stress (e.g. from as yet unexperienced environmental
  stressors) and in local adaptation.

  In the present project, differential gene-expression in
  \textit{A. crassus} populations was assessed using next generation
  sequencing on the 454 and Illumina platforms and genetic components
  of differences were isolated in cross-inoculation experiments with
  both Asian and European host-species and parasite populations.

  Several proteases were shown to be under positive selection on the
  sequence level, highlighting this group of enzymes as possible
  targets of an immune-attack on \textit{A. crassus}. On the
  gene-expression level, the extent of heritable change was large in
  comparison to the effect of modification in different
  host-environments. Mitochondrially encoded subunits of the
  respiratory chain and other genes connected to aerobic respiration
  showed divergent expression patterns in European vs. Asian parasite
  populations; cuticle collagen genes showed ``adapted'' patterns of
  expression in present day sympatric host-parasites pairs.
  
  These results identified gene-expression phenotypes, confirming the
  divergence of European \textit{A. crassus} populations. Such
  phenotypes will be more accessible to popluation-genetic analysis
  invetigating selection than complex life history traits.

\end{abstractslong}
%\end{abstractlongs}


% ---------------------------------------------------------------------- 

%%% Local Variables: ***
%%% mode:latex ***
%%% TeX-master: "../thesis.tex"  ***
%%% tex-main-file: "../thesis.tex" ***
%%% End: ***
     
% \end{abstractseparate}


%: ----------------------- tie in front matter ------------------------

\frontmatter
% Thesis Dedictation ---------------------------------------------------

\begin{dedication} %this creates the heading for the dedication page

  Natur! [...] Es ist ein ewiges Leben, Werden und Bewegen in ihr, und
  doch rückt sie nicht weiter. Sie verwandelt sich ewig, und ist kein
  Moment Stillestehen in ihr. Fürs Bleiben hat sie keinen Begriff, und
  ihrem Fluch hat sie ans Stillestehen gehängt. Sie ist fest. Ihr
  Tritt ist gemessen, ihre Ausnahmen selten, ihre Gesetze unwandelbar.


  NATURE! [...] Incessant life, development, and movement are in her,
  but she advances not. She changes for ever and ever, and rests not a
  moment. Quietude is inconceivable to her, and she has laid her curse
  upon rest. She is firm. Her steps are measured, her exceptions rare,
  her laws unchangeable.

  J. W. GOETHE translated to English by T. H. HUXLEY for the first
  issue of nature magazine.

  \vspace{6cm}

  To my grandmother Ruth my brother Roman and my wife Silvia

\end{dedication}

% ----------------------------------------------------------------------
% Thesis Acknowledgements ------------------------------------------------


%\begin{acknowledgementslong} %uncommenting this line, gives a different acknowledgements heading
\begin{acknowledgements}      %this creates the heading for the acknowlegments


  I would like to acknowledge the following persons for their
  tremendous help realizing this project (in alphabetical order):

  Mark Blaxter, Stephen Bridgett, Timothee Cezard, Yun-San Han, Sujai
  Kumar, Dominik Laetsch, Jenna Mann, Anna Montazam, Horst
  Taraschewski, Marian Thomson, Urmi Trivedi, Urszula Weclawski,
  Nicola Wrobel and the thousands of unnamed individuals who have
  coded for free software and open source projects used in my
  research.

\end{acknowledgements}
%\end{acknowledgmentslong}

% ------------------------------------------------------------------------





%: ----------------------- contents ------------------------

\setcounter{secnumdepth}{3} % organisational level that receives a numbers
\setcounter{tocdepth}{3}    % print table of contents for level 3
\tableofcontents            % print the table of contents
% levels are: 0 - chapter, 1 - section, 2 - subsection, 3 - subsection


%: ----------------------- list of figures/tables ------------------------

\listoffigures	% print list of figures

\listoftables  % print list of tables


%: ----------------------- glossary ------------------------

% Tie in external source file for definitions: /0_frontmatter/glossary.tex
% Glossary entries can also be defined in the main text. See glossary.tex
% this file is called up by thesis.tex
% content in this file will be fed into the main document

% Glossary entries are defined with the command \nomenclature{1}{2}
% 1 = Entry name, e.g. abbreviation; 2 = Explanation
% You can place all explanations in this separate file or declare them in the middle of the text. Either way they will be collected in the glossary.

% required to print nomenclature name to page header
\markboth{\MakeUppercase{\nomname}}{\MakeUppercase{\nomname}}


% ----------------------- contents from here ------------------------

%  
\nomenclature{SNP}{Single Nucleotide Polymorphism; variation occurring
in a single nucleotide between two closely related homlogous
sequences. Leading to for example to allelic differences within a
population or even the homologous chromosomes in an individual}


\nomenclature{ORF}{Open Reading Frame; a region in a DNA-sequence
begining with a start-codon and not containing a stop-condon. For
example a region within a processed mRNA transcript being transcribed
into a protein}

\nomenclature{dpi}{Days post infection; In infection experiments, a
point in time given in days after an individual has been infected}

\nomenclature{DNA}{Desoxy Ribonucleic Acid; a chemical molecule
  bearing the heritable genetic information in all life on earth}


 

\begin{multicols}{2} % \begin{multicols}{#columns}[header text][space]
\begin{footnotesize} % scriptsize(7) < footnotesize(8) < small (9) < normal (10)

\printnomenclature[1.5cm] % [] = distance between entry and description
\label{nom} % target name for links to glossary

\end{footnotesize}
\end{multicols}



%: --------------------------------------------------------------
%:                  MAIN DOCUMENT SECTION
% --------------------------------------------------------------

% the main text starts here with the introduction, 1st chapter,...
\mainmatter

\renewcommand{\chaptername}{} % uncomment to print only "1" not "Chapter 1"


%: ----------------------- subdocuments ------------------------

% Parts of the thesis are included below. Rename the files as required.
% But take care that the paths match. You can also change the order of appearance by moving the include commands.


% this file is called up by thesis.tex
% content in this file will be fed into the main document

%: ----------------------- introduction file header -----------------------
\chapter{Introduction}

% the code below specifies where the figures are stored
\ifpdf
    \graphicspath{{1_introduction/figures/PNG/}{1_introduction/figures/PDF/}{1_introduction/figures/}}
\else
    \graphicspath{{1_introduction/figures/EPS/}{1_introduction/figures/}}
\fi

% ----------------------------------------------------------------------
%: ----------------------- introduction content ----------------------- 
% ----------------------------------------------------------------------



%: ----------------------- HELP: latex document organisation
% the commands below help you to subdivide and organise your thesis
%    \chapter{}       = level 1, top level
%    \section{}       = level 2
%    \subsection{}    = level 3
%    \subsubsection{} = level 4
% note that everything after the percentage sign is hidden from output



\section{The study organism: \textit{Anguillicola crassus}} 

\subsection{Ecological significance} 

\textit{Anguillicola crassus} Kuwahara, Niimi and Ithakagi 1974
\cite{kuwahara_Niimi_Itagaki_1974, moravec_anguillicoloides} is a
swimbladder nematode naturally parasitizing the Japanese eel
(\textit{Anguilla japonica}) indigenous to East-Asia. After a single
introduction \cite{wielgoss_population_2008} to Germany in the early
1980s \textit{A. crassus} has colonized almost all populations of the
European eel (\textit{Anguilla anguilla})
\cite{kirk_impact_2003}. Since the 1990s populations of the American
eel (\textit{Anguilla rostrata}) have been colonized as novel hosts
\cite{fries_notes:_1996,barse_exotic_1999, barse_swimbladder_2001} and
finally it has been detected in three indigenous \textit{Anguilla}
species on the island of Reunion near Madagascar
\cite{sasal_parasite_2008}.\



In Asia, as well as in the introduced ranges, copepods and ostracods
serve as intermediate hosts of \textit{A. crassus}
\cite{moravec_first_2005}, in which L2 larvae develop to L3 larvae,
infective to the final host. Once ingested by an eel they migrate
through the intestinal wall and via the body cavity into the
swimbladder wall \cite{haenen_effects_1996}, i.a. using a trypsin-like
proteinase\cite{polzer_identification_1993}. In the swimbladder wall
L3 larvae hatch to L4 larvae. After a final moult from L4 to preadult
the parasites inhabit the lumen of the swimbladder, where they
eventually mate. Eggs containing L2 larvae are released via the ductus
pneumaticus into the eels gut and finally into the
water\cite{de_charleroy_life_1990}.\

Within the novel range and hosts, conspicuously elevated prevalences
and intensities of infection occur (reviewed in
\cite{kirk_impact_2003} and \cite{taraschewski_hosts_2007}). These
differences in abundance of \textit{A. crassus} in East Asia compared
to Europe are commonly attributed to the different host-parasite
relations in the final eel host permitting a differential survival of
the larval and the adult parasites
\cite{knopf_differences_2004}. Recently, data from experimental
infections of European eels with \textit{A. crassus} have been
published \cite{fazio_regulation_2008}. They show that the parasite
undergoes (under experimental conditions) a density-dependent
regulation keeping the number of worms within a certain range.\
  

\subsection{Evolutionary significance}

\subsubsection{Divergence of \textit{A. crassus} populations}

\figuremacro{ula_stages.png}{Differences in developmental speed}{data
  courtesy of Urszula Weclawski
  %% \href{http://someurl}{urlanchor<}.
  %% \figuremacroW{ula_stages.png}{Title}{Caption}{0.8}
  % variation of the above macro with a width setting
}

\subsubsection{Interest in \textit{A. crassus based on its
    phylogenetic position in the phylum nematoda}}



\subsection{Functional insights from other nematodes used to formulate
  hypotheses for \textit{A.crassus}}


\section{Advances in sequencing technology enabeling this study}

Recent advances in DNA-sequencing


\figuremacro{sequencing_costs.png}{Falling seqeuncing costs}{Falling into bottomless, Data
  provided by \href{http://www.genome.gov/sequencingcosts/}{National
    Human Genome Research Institute, NHGRI}.}



\subsection{Pyro-sequencing}


\subsection{Illumina-Solexa sequencing}



\begin{table}[htdp]
\centering
\begin{tabular}{ccc} % ccc means 3 columns, all centered; alternatives are l, r

{\bf Gene} & {\bf GeneID} & {\bf Length} \\ 
% & denotes the end of a cell/column, \\ changes to next table row
\hline % draws a line under the column headers

human latexin & 1234 & 14.9 kbps \\
mouse latexin & 2345 & 10.1 kbps \\
rat latexin   & 3456 & 9.6 kbps \\
% Watch out. Every line must have 3 columns = 2x &. 
% Otherwise you will get an error.

\end{tabular}
\caption[title of table]{\textbf{title of table} - Overview of latexin genes.}
% You only need to write the title twice if you don't want it to appear in bold in the list of tables.
\label{latexin_genes} % label for cross-links with \ref{latexin_genes}
\end{table}

% There you go. You already know the most important things.


% ----------------------------------------------------------------------

%%% Local Variables: ***
%%% mode:latex ***
%%% TeX-master: "../thesis.tex"  ***
%%% End: ***
     	% background information
% this file is called up by thesis.tex
% content in this file will be fed into the main document

\chapter{Aims of the project} % top level followed by section, subsection

% ----------------------- contents from here ------------------------

\section{Final aim}

In a reciprocal transplant experiment genetic components of the
differences in 







\section{Preliminary aims}

In order to investigate transcriptomic response to environmental
stimuli, the responding unit, the transcripts have to be established
first. As extremely short reads providing ultra high thoroughput are
hard to assemble \textit{de-novo}, a reference was created first using
454 pyrosequencing technology providing longer read-length.



% ---------------------------------------------------------------------------
% ----------------------- end of thesis sub-document ------------------------
% ---------------------------------------------------------------------------

%%% Local Variables: ***
%%% mode:latex ***
%%% TeX-master: "../thesis.tex"  ***
%%% tex-main-file: "../thesis.tex" ***
%%% End: ***
     						% aims of the project

% this file is called up by thesis.tex
% content in this file will be fed into the main document

%: ----------------------- name of chapter  -------------------------
\chapter{Pilot sequencing (Sanger method)} % top level followed by section, subsection
\label{pilot-seq}

%: ----------------------- paths to graphics ------------------------

% change according to folder and file names
\ifpdf
    \graphicspath{{3_pilot/figures/PNG/}{3_pilot/figures/PDF/}{3_pilot/figures/}}
\else
    \graphicspath{{3_pilot/figures/EPS/}{3_pilot/figures/}}
\fi

%: ----------------------- contents from here ------------------------

\subsection{Overview}

This chapter reports a small pilot-project investigating the
RNA-extraction and cDNA preparation in preparation for high-throughput
transcriptome sequencing of the swimbladder nematode
\textit{A. crassus}. I generated expressed sequence tags (ESTs) using
traditional Sanger-technology and conducted a first assessment of the
sequence diversity expected in deeper sequencing. Especially the
expected coverage of unwanted rRNA and host-derived sequences was
investigated.

In total 945 reads from adult \textit{A. crassus} (5 libraries from 4
cDNA preparations, including 541 sequences generated by students in a
laboratory course) and 288 reads from liver-tissue of the host species
\textit{An. japonica} (3 libraries from 3 cDNA preparations) were
sequenced.

\subsection{Initial quality screening}
\label{qual-pil}

The initial quality screening revealed a high number of sequences that
had to be discarded due to failed sequencing reactions (sequences
being too short after quality trimming by \texttt{trace2seq}) in the
library prepared by students. For sequences of \textit{An. japonica}
and the other libraries from \textit{A. crassus}, failed sequencing
reactions were less common.

In the next screening-step for \textit{A. crassus} 125 (13.23\%) and
for \textit{An. japonica} 64 (22.22\%) of the sequences were
excluded because of homopolymer-runs considered to be artificial. This
resulted in 452 of the nematode and 195 of the host reads being
regarded of sufficient quality for further processing after
base-calling and quality screening.

\subsection{rRNA screening}
\label{rRNA-pil}

The further screening of sequences revealed a high abundance of rRNA
(see Figure \ref{rRNA-con}) ranging from 71.67\% to 91.67\% of the
obtained sequences.  High abundances of rRNA were also found in the
libraries from host liver tissue (see table \ref{tab:num}), ranging
from 71.67\% to 77.42\%. This contamination in libraries from both
species was mainly responsible for a low number of sequences being of
sufficient quality for submission to NCBI-dbEST. At this point for the
\textit{An. japonica}-dataset, 36 sequences were submitted to
NCBI-dbEST under the Library Name ``\textit{Anguilla japonica} liver''
and were assigned the accession LIBEST\_027503.

\figuremacro{rRNA-con}{Proportion of rRNA in different libraries for
  \textit{A. crassus} and \textit{An. japonica}}{rRNA abundance as
  proportion of the raw sequencing-reads (rRNA from total) and as
  proportion of the reads after quality screening (rRNA from
  good). Libraries starting with ``Ac\_'' are from \textit{A. crassus},
  libraries starting with ``Aj\_'' are from \textit{An. japonica}.}

% latex table generated in R 2.13.0 by xtable 1.5-6 package
% Mon Sep 19 15:44:26 2011
\small\begin{table}[ht]
\begin{center}
\begin{tabular}{rrrrrr}
  \hline
 & short & poly & rRNA & fishpep & good \\ 
  \hline
Ac\_197F(n=96) &   4 &  17 &  58 &   1 &  16 \\ 
  Ac\_106F(n=96) &  25 &   9 &  48 &   0 &  14 \\ 
  Ac\_M175(n=116) &  30 &  19 &  41 &   3 &  23 \\ 
  Ac\_FM(n=96) &  12 &  29 &  34 &   1 &  20 \\ 
  Ac\_EH1(n=541) & 297 &  51 & 143 &   8 &  42 \\ 
  Ac\_total(n=945) & 368 & 125 & 324 &  13 & 115 \\ 
  Aj\_Li1(n=96) &  10 &  23 &  50 &  &  13 \\ 
  Aj\_Li2(n=96) &  10 &  26 &  43 &  &  17 \\ 
  Aj\_Li3(n=96) &   9 &  15 &  66 &  &   6 \\ 
  Aj\_total(n=288) &  29 &  64 & 159 &  &  36 \\ 
   \hline
\end{tabular}
\caption[Screening statistics for pilot sequencing] {\textbf{Screening
    statistics for pilot sequencing} - Number of ESTs discarded at
  each screening-step for single libraries and totals for
  species. Short, sequence to short in \texttt{trace2seq}; poly,
  sequences with artificial homopolymer-runs from poly-A tails; rRNA,
  with hits to rRNA databases; fishpep, with better hits to
  host-protein-databases than to nematode protein databases; good,
  sequences regarded ``valid'' after all screening steps. Note that
  the 13 sequences in the \textit{A. crassus}-dataset, for which
  fish-origin was inferred, were still submitted to NCBI-dbEST.}

\label{tab:num}
\end{center}
\end{table}
\normalsize

\subsection{Screening for host-contamination}
\label{host-pil}

For the \textit{A. crassus}-dataset screening for host-sequences at
this stage was regarded necessary based on the notion that a large
proportion of the tissue prepared in RNA extraction consisted of
eel-blood inside the gut of the worms (see also Figure
\ref{worm_diff}). Additionally, a bimodal distribution of GC-content
in the \textit{A. crassus}-dataset was observed with one of the modes
consistent with the mean GC-content of the ESTs from the Japanese eel.

Comparison of \texttt{Blast}- results for these sequences versus
nempep4 and a fishprotein-database (derived from NCBI non-redundant),
showed that 13 sequences were more likely to originate from host
contamination than from \textit{A. crassus}. These 13 sequences in the
\textit{A. crassus} data-set were submitted to NCBI-dbEST with a
comment that host origin had been inferred. This reduced the dataset
essentially to 115 ESTs. However, these 13 ESTs are still accessible
through the same library name ``Adult \textit{Anguillicola crassus}''
and library-identifier LIBEST\_027505 and are taxonomically attributed
to \textit{A. crassus} on NCBI-dbEST.

\figuremacro{GC-pil}{GC-content of sequences from \textit{An. japonica}
  and \textit{A. crassus}}{The Japanese eel has a slightly higher
  GC-content than the parasite: This sequence characteristic is useful
  for separation of sequences from the host-parasite interface, note
  the higher GC-content of the sequences from \textit{A. crassus}, for
  which host origin was inferred from similarity searches (red line
  labeled \textit{A. crassus/An .japonica}).}

After screening of host-sequences the GC-content of
\textit{A. crassus} ESTs had a unimodal distribution (see Figure
\ref{GC-pil}). \textit{A. crassus} had a lower mean GC-content $(37.32
\pm 8.36$ mean $\pm$ sd) than \textit{An. japonica} $(45.79 \pm
8.36$ mean $\pm$ sd; two-sided t-test $p<0.001$). The distribution of
the GC-contents for sequences, for which host-origin was inferred was
in agreement with the GC-distribution for host sequences.

\texttt{Blast}-annotations obtained (by similarity searches against
NCBI-nr, bit-score threshold of 55) for the sequences of putative host
origin were also largely in agreement with the expectations for
eel-blood: one sequence could be identified being highly similar to
``hemoglobin anodic subunit'' from the European eel. Others were
annotated with best hits to highly expressed housekeeping genes from
fish or vertebrates (see table \ref{tab:hostan}). Two sequences in the
set had lower similarities only to proteins predicted from
genome-sequences of chordates, and one sequence of the 13 lacked any
similarity to NCBI-nr above the threshold of 55 bits.

115 of the submitted sequences for ``Adult \textit{Anguillicola
  crassus}'' (LIBEST\_027505) were regarded as ``valid'', i.e. not
clearly of host origin.

However, two ESTs (Ac\_EH1f\_01D10 and Ac\_EH1r\_01D10; forward and
reverse read of the same clone) were annotated with
``ref$|$ZP\_05032178.1$|$; exopolysaccharide synthesis, ExoD
superfamily'' from \textit{Brevundimonas} sp. BAL3. The family
Caulobacteraceae, comprises bacteria living in freshwater and
sequences are probably derived from a commensal, symbiont or pathogen
of eels or swimbladder-nematodes. These off-target data were left in
the submission file.

For 66 (58.4\%) of the remaining 113 ESTs annotations were obtained
from orthologous sequences. All of these orthologous sequences were
from other species in the phylum nematoda.

% Sun Oct  2 07:46:25 2011
\begin{sidewaystable}[ht]
\begin{center}
\begin{tabular}{lrp{3.5cm}lrr}
  \hline
  sequence & hit identifier & hit description &species & bit-score & e-value \\
  \hline
  Ac\_EH1f\_005B07 & gb$|$AAQ97992.1$|$ & cyclin G1 & \textit{Danio rerio} & 67.0 & 9e-10 \\ 
  Ac\_EH1f\_01A02 & gb$|$ACO10003.1$|$ &  Nicotinamide riboside kinase 2 & \textit{Osmerus mordax} &  333 & 1e-89 \\ 
  Ac\_EH1f\_01C10 & gb$|$ADF80517.1$|$ & ferritin M subunit & \textit{Sciaenops ocellatus}  &  328 & 5e-88 \\ 
  Ac\_EH1r\_004A04 & ref$|$XP\_003340320.1$|$ & cytoplasmic 1-like actin & \textit{Monodelphis domestica}  &  102 &3e-20 \\ 
  Ac\_EH1r\_005B07 & gb$|$ABN80454.1$|$ & cyclin G1  &\textit{Poecilia reticulata}  & 90.5 & 8e-17 \\ 
  Ac\_EH1r\_009C03 & ref$|$NP\_001122208.1$|$ & THAP domain containing protein 4 & \textit{Danio rerio}  &  176 & 1e-42 \\ 
  Ac\_EH1r\_01A07 & sp$|$P80946.1$|$ & Hemoglobin anodic subunit beta &\textit{Anguilla anguilla} & 283 &1e-74 \\ 
  Ac\_FMf\_08F03 & ref$|$XP\_003226802.1$|$ & cohesin subunit SA-2-like isoform 2 & \textit{Anolis carolinensis}  &  219 &8e-56 \\ 
  Ac\_M175\_01H02 & emb$|$CAQ87569.1$|$ & NKEF-B protein & \textit{Plecoglossus altivelis}  &  365 &3e-99 \\ 
  Ac\_197Ff\_01E04 & ref$|$XP\_002121150.1$|$ & CUB and sushi domain-containing protein 3 & \textit{Ciona intestinalis} & 80.5 & 2e-13 \\ 
  Ac\_EH1f\_01D07  & ref$|$XP\_002606965.1$|$ & hypothetical protein & \textit{Branchiostoma floridae} & 82.8 &3e-14 \\ 
  Ac\_M175\_01B06 & ref$|$XP\_422710.2$|$ & hypothetical protein & \textit{Gallus gallus} &  123 &1e-26 \\
  \hline
\end{tabular}
\end{center}
\caption[Annotation of putative host-derived sequences in the
\textit{A. crassus}-dataset]{\textbf{Annotation of putative
    host-derived sequences in the \textit{A. crassus}-dataset} -
  Sequences excluded because of inferred host-origin comparing
  similarity to nematode- and fish-proteins. The annotation obtained
  against NCBI-nr are in agreement with this inference of host origin,
  as only best hits to vertebrate proteins are found.}
\label{tab:hostan}
\end{sidewaystable}

%%% Local Variables: ***
%%% mode:latex ***
%%% TeX-master: "../thesis.tex"  ***
%%% tex-main-file: "../thesis.tex" ***
%%% End: ***
     
			
% this file is called up by thesis.tex
% content in this file will be fed into the main document

%: ----------------------- name of chapter  -------------------------
\chapter{Pyrosequencing of the \textit{A. crassus}
  transcriptome} % top level followed by section, subsection


%: ----------------------- paths to graphics ------------------------

% change according to folder and file names
\ifpdf
    \graphicspath{{X/figures/PNG/}{X/figures/PDF/}{X/figures/}}
\else
    \graphicspath{{X/figures/EPS/}{X/figures/}}
\fi

%: ----------------------- contents from here ------------------------

454 







%%% Local Variables: ***
%%% mode:latex ***
%%% TeX-master: "../thesis.tex"  ***
%%% tex-main-file: "../thesis.tex" ***
%%% End: ***

	
% this file is called up by thesis.tex
% content in this file will be fed into the main document

%: ----------------------- name of chapter  -------------------------
\chapter{NlaIII-tag sequencing
  (Super-SAGE)} % top level followed by section, subsection


%: ----------------------- paths to graphics ------------------------

% change according to folder and file names
\ifpdf
    \graphicspath{{X/figures/PNG/}{X/figures/PDF/}{X/figures/}}
\else
    \graphicspath{{X/figures/EPS/}{X/figures/}}
\fi

%: ----------------------- contents from here ------------------------

\section{Comparison with pyrosequencing-data}




%%% Local Variables: ***
%%% mode:latex ***
%%% TeX-master: "../thesis.tex"  ***
%%% tex-main-file: "../thesis.tex" ***
%%% End: ***

% this file is called up by thesis.tex
% content in this file will be fed into the main document

%: ----------------------- name of chapter  -------------------------
\chapter{Transcriptomic divergence in a common garden experiment}
% top level followed by section, subsection
\label{cha:tra-diff}

%: ----------------------- paths to graphics ------------------------

% change according to folder and file names
\ifpdf
    \graphicspath{{6_rna_seq/figures/PNG/}{6_rna_seq/figures/PDF/}{6_rna_seq/figures/}}
\else
    \graphicspath{{6_rna_seq/figures/EPS/}{6_rna_seq/figures/}}
\fi


%: ----------------------- contents from here ------------------------


\section{Infection experiments}

Dissection of eels 55-57 dpi showed higher recovery of European worms
in \textit{An. anguilla} and higher recovery of Taiwanes worms in
\textit{An. japonica}, compared to the other parasite populations. In
other words, in host-parasite combinations of matching origin,
parastites performed better.

\figuremacro{stages_plot}{Recovery of worms in cross-infection
  experiment}{Mean numbers of worms recovered after 55-57 dpi for
  sample sizes given as n=x. Error-bars indicate the standard error
  (s.e.) of the mean. Stages are L3-larvae (l3), L4-larvae (l4), adult
  females (adult.f) and adult males (adult.m).}

In host-species/parastite-population pairs found in nature roughly
eight or nine adult worms could be recoverd per eel. In the
transplanted host/parasite combinations only two or three adult worms
were recovered on average (see figure \ref{stages_plot}). In
\textit{An. anguilla} no differences in the recovery of larval stages
was recorded. In \textit{An. japonica} however roughly two individals
more were recorded from both larval stages in the host/parasite
combination found in nature.

These differences are highly significant especially for adult worms
(see table \ref{tab:ad-sig}) and are interpretable as a sign of local
adaptation, as adult survival and recovery can be regarded a fitness
component.

Recovery as a proportion of the 50 larvae eels were inocculated with,
was thus roughly 30\% for the adapted pairs compared to only roughly
10\% in non-adapted host-parasite pairs.

\begin{table}[ht]
\begin{center}
\begin{tabular}{rrrrr}
  \hline
 & Estimate & Std. Error & t value & Pr($>$$|$t$|$) \\ 
  \hline
  (Intercept) & 9.5000 & 1.1109 & 8.55 & 0.0000 \\ 
  host.spec.AJ & -8.0789 & 1.7472 & -4.62 & 0.0000 \\ 
  worm.pop.T & -5.2222 & 1.3689 & -3.81 & 0.0002 \\ 
  host.spec.AJ:worm.pop.T & 11.7345 & 2.2010 & 5.33 & 0.0000 \\ 
   \hline
\end{tabular}
\caption[Linear model for recovery]{Linear model for recovery of adult
  worms. The estimate gives the mean of the distribution of adult
  worms for the factor values in the rows. The intercept is set to Aa
  R (\textit{An. anguilla} and the European population) further rows
  give variations for each factor. Std. Error is the standard error of
  this value. Additonally the probabilty of a t-value as small or
  smaller then the oberserved t-value are given. The signature of
  local adaptation is visible in the highly significant interaction
  term.}
\label{tab:ad-sig}
\end{center}
\end{table}
\afterpage{\clearpage}

\section{Sample preparation and sequencing }

Three biological replicates were obtained from each of the two worm
population in each of the two eel-host for each of the two sexes of
worms. This resulted in a total 24 RNA-extractions prepared for
sequencing: 3 individual female worms from each experimental group
were choosen randomly, to give in total twelve females. Additionally
from three individual male worms, and from 9 pools of male worms RNA
was extracted (see table \ref{tab:lib-prep}). Pools consisted of worms
from one infected eel individual each. All worms or worm-pools were
derived form infections of different eel individuals, with one small
exception from this form of statistical independence: From eel AJ/R3 a
male worms as well as a female worm had to be prepared. It was
impossible to extract enough RNA from all but the biggest male worms
especially of the Japanese eel/European worm combination, leaving no
other choice. Because of the small size of male worms it was generally
not possible to randomly choose individuals. Preparation of sufficient
amounts of RNA was only achieved in pools of the biggest
individuals. All male worms were thus choosen for preparation based on
large size, even when pools of worms were used.

% latex table generated in R 2.13.0 by xtable 1.5-6 package
% Tue Nov 29 10:05:46 2011
\begin{table}[h]
\begin{center}
\begin{tabular}{llllrrr}
  \hline
label & sex & host & population & intensity & worms in prep & conc in prep \\ 
  \hline
AA/T20F & female & AA & Taiwan & 1 & 1 & 5.60 \\ 
  AA/T12F & female & AA & Taiwan & 14 & 1 & 6.80 \\ 
  AA/T45F & female & AA & Taiwan & 5 & 1 & 8.00 \\ 
  AA/T24M & male & AA & Taiwan & 6 & 3 & 4.80 \\ 
  AA/T42M & male & AA & Taiwan & 11 & 1 & 5.60 \\ 
  AA/T3M & male & AA & Taiwan & 5 & 4 & 4.88 \\ 
  AA/R18F & female & AA & Europe & 4 & 1 & 4.80 \\ 
  AA/R28F & female & AA & Europe & 10 & 1 & 5.20 \\ 
  AA/R8F & female & AA & Europe & 27 & 1 & 5.20 \\ 
  AA/R16M & male & AA & Europe & 10 & 4 & 5.20 \\ 
  AA/R11M & male & AA & Europe & 25 & 14 & 6.40 \\ 
  AA/R2M & male & AA & Europe & 10 & 4 & 6.60 \\ 
  AJ/T8F & female & AJ & Taiwan & 10 & 1 & 5.91 \\ 
  AJ/T5F & female & AJ & Taiwan & 2 & 1 & 4.80 \\ 
  AJ/T26F & female & AJ & Taiwan & 2 & 1 & 2.40 \\ 
  AJ/T25M & male & AJ & Taiwan & 24 & 5 & 4.05 \\ 
  AJ/T19M & male & AJ & Taiwan & 24 & 7 & 3.50 \\ 
  AJ/T20M & male & AJ & Taiwan & 20 & 8 & 3.80 \\ 
  AJ/R1F & female & AJ & Europe & 3 & 1 & 5.92 \\ 
  AJ/R3F & female & AJ & Europe & 3 & 1 & 6.90 \\ 
  AJ/R5F & female & AJ & Europe & 10 & 1 & 4.04 \\ 
  AJ/R1M & male & AJ & Europe & 3 & 1 & 2.50 \\ 
  AJ/R3M & male & AJ & Europe & 3 & 2 & 2.60 \\ 
  AJ/R5M & male & AJ & Europe & 10 & 1 & 2.23 \\ 
   \hline
\end{tabular}
\caption[Summary of RNA preparation]{\textbf{A summary of all 24
    samples prepared for RNA-seq -} The label of the RNA preparation
  follows a convention based on the eel species (host; first two
  letter of label, AA for \textit{An. anguilla} AJ for
  \textit{An. japonica}), worm population (population - R for
  European, T for Taiwanese; letter after the slash) and sex of
  worm(s) in preparation (F for female, M for male; last letter in
  label). Additionally the intensity of infection (number of adult
  worms found in the infected eel; intensity) and the number of worms
  pooled in the preparation (only male worms are pooled for RNA
  extraction, individual female worms were used). Finally
  RNA-concentration (conc in prep) in the preparation is given in
  $\mu$g per ml}
\label{tab:lib-prep}
\end{center}
\end{table}

\afterpage{\clearpage}

Sequencing was performed in three multiplexed pools of eight libraries
each. The samples were partitioned into these pools spreading
replicates over lanes in a blocking design to further guarantee
statistical independence from sequencing-lane effects. Each pool of
eight was sequenced on two lanes, giving in total six lanes of data
and two technical replicates for each library. Sequencing resulted in
a total of 263.668.952 raw sequencing read-pairs.

\section{Examination of data-quality}

Reads were mapped against the fullest assembly (see
\ref{sec:final-full-assembly}) using \texttt{BWA}
\cite{pmid20080505}. Of the 263.668.952 raw read-pairs 173.602.387
mapped uniquely to the assembly and were counted on a per-library
base.

Read counts for one library were originally obtained from two
different technical replicate lanes of a flow-cell. The technical
replicate of read-counts demonstrated very low differences as inferred
from a clustering analysis using variance stabilized data and
euclidian distances as implemented in \texttt{DESeq}
\cite{pmid20979621} (see \ref{heat_tech_rep}).

\figuremacroW{heat_all}{Distances between RNA-seq read-count for
  different samples}{Euclidean distance (square distance between the
  two count vectors) for variance stabilized read-counts for all
  libraries including technical replicates; Red indicates low distance
  (high similarity), blue heigh distance (low similarity). a) Data
  before screening and summation of technical replicates. All
  technical replicates are clusterd very closely, the distance between
  an outlier female sample (AJ\_T26F) is high. b) Same illustration
  after summation of technical replicates and screening. Distance
  between outlier-sample and other female samples is reduced.}{0.5}

158.232.523 read-pairs were left after removal of hits to contigs for
which non-\textit{A. crassus} origin had been inferred in the
curration of the 454-transcriptome assembly.

% latex table generated in R 2.13.0 by xtable 1.5-6 package
% Tue Nov 29 14:12:43 2011
\begin{table}[h]
\begin{center}
\begin{tabular}{llrrr}
  \hline
library & raw.reads & raw.mapped & tax.mapped & screened \\ 
  \hline
AA\_R11M & 11986442 & 8628520 & 7868814 & 6889551 \\ 
  AA\_R16M & 10810349 & 6858585 & 6217540 & 5276284 \\ 
  AA\_R18F & 9227615 & 6552527 & 5933235 & 5200958 \\ 
  AA\_R28F & 10135670 & 6665381 & 6005399 & 5171806 \\ 
  AA\_R2M & 12469746 & 7628428 & 6929651 & 5906422 \\ 
  AA\_R8F & 15270570 & 11527867 & 10758535 & 9453468 \\ 
  AA\_T12F & 11299438 & 7842479 & 7195621 & 6332396 \\ 
  AA\_T20F & 11740839 & 7744179 & 7114349 & 6323422 \\ 
  AA\_T24M & 8552723 & 5254194 & 4662053 & 3969305 \\ 
  AA\_T3M & 11031751 & 6460836 & 5800042 & 4993726 \\ 
  AA\_T42M & 11573501 & 7567845 & 6787375 & 5694801 \\ 
  AA\_T45F & 10646847 & 7714472 & 7173709 & 6283585 \\ 
  AJ\_R1F & 9855005 & 6400558 & 5890748 & 5167912 \\ 
  AJ\_R1M & 10211903 & 5851063 & 5313544 & 4506254 \\ 
  AJ\_R3F & 9897937 & 6425201 & 5948079 & 5124077 \\ 
  AJ\_R3M & 8775211 & 4562324 & 4073621 & 3422526 \\ 
  AJ\_R5F & 11949105 & 8442537 & 7830247 & 6882280 \\ 
  AJ\_R5M & 11231532 & 7504494 & 6772010 & 5913016 \\ 
  AJ\_T19M & 9195576 & 4798404 & 4293123 & 3635843 \\ 
  AJ\_T20M & 10862591 & 6880937 & 6251674 & 5280529 \\ 
  AJ\_T25M & 11195315 & 7162880 & 6480185 & 5645097 \\ 
  AJ\_T26F & 11195335 & 7439917 & 6641973 & 6031374 \\ 
  AJ\_T5F & 10357569 & 7413685 & 6794507 & 6007930 \\ 
  AJ\_T8F & 14196382 & 10275074 & 9496489 & 8364594 \\ 
   \hline
\end{tabular}
\caption[Mapping Summary]{\textbf{Maping summarized for all 24
    libraries -} Rows indicate different libraries (worms or
  worm-pools as indicated in \ref{tab:lib-prep}) raw.reads gives the
  number of read-paires sequenced, raw.mapped the number of reads
  mapping uniquely with their best hit, tax.mapped the number of reads
  after substraction of reads to putatively eel-host derived contigs
  and screened after substraction of all reads mapping not to the
  higCA-derived assembly or to contigs with overall counts less than
  32.}
\label{tab:read-clean}
\end{center}
\end{table}

\afterpage{\clearpage}

After another screening for spurious read-counts to low covered
transcripts and to transcripts of low relieability (lowCA in the
454-assembly; see \ref{sec:final-full-assembly}) 137.477.156
read-pairs were left for further analysis. Distribution of these
read-pairs over libraries showd roughly three fold differences, with a
mean of 5.728.215 reads and a range from 3.422.526 read-pais for
library AJ\_R3M to 9.453.468 read-pairs for library AA\_R8F (see
\ref{tab:read-clean}).

\figuremacro{mds}{Principle coordinate plot for expression inRNA-seq
  libraries}{Distance between sample-pairs is the root-mean-square
  deviation (Euclidean distance) for the most differentially expressed
  (DE) genes. Distances can be interpreted as the log2-fold-change of
  the genes with the biggest changes, i.e. the log2-fold-change for
  the genes that distinguish the samples.}

These reads mapped to 7520 contigs from our 454 assembly, making them
the basis for all further inverstigations. Analysis of between-sample
distance confirmed the library clustering. Sex of the worms defined
the overall distances between libraries, host- or
population-differences were not visible in an overall effect in the
top differentially expressed (DE) genes (see figure \ref{mds}).

\section{Orthologous-screened expression differences}

For the 7520 contigs with analysed exprssion values 4382
\textit{C. elegans}-orthologs and 4292 \textit{B. malayi}-orthologs
with read more than 32 read-counts over all libraries were determined
based on the annotation of the assembly (see \ref{454-annot}). This
resulted in 3596 contigs with measured expression also having a
measurement for a corresponding orhtolog (or group of orthologs) in
both model-species and thus being available for analysis.

For all further evaluations the congruence of the basic contig-based
statistics with orthologous-derived statistics is considered.

\section{Expression differences in general linear models}

Generalized linear models (GLMs) were used as implemented in the
R-package \texttt{edgeR}. Using these models we obtained 2588 contigs
(34\% of total) DE between male and female worms at a false discovery
rate (FDR) of 5\%. 1101 (31\% of total orthologous available) of these
contigs of were confirmed by contigs in the orthologous
evaluation. 1425 (556 orthologous confirmed; OC) of these were
upregulated in male worms 1163 (545 OC) upregulated in female worms.

At the same threshold 55 contigs (0.7\% of total; 9, 0.25\% OC) showed
significant differential response to the host-species. 38 (5 OC) were
upregulated upregulated in \textit{An. japonica}, 17 (4 OC)
upregulated in \textit{An. anguilla}.

68 contigs (0.9\% of total; 15, 0.42\% OC) showed differences
according to the population of the worm. 39 (11 OC) of these were
upregulated in the Taiwanese population, 29 (4 OC) in the European
population.

An important observation in these models is the prevalence of
co-occuring significance of simple main effects. Expression changes
overlapping for two main effects mean a significant difference in
expression according to both factors. These differences are in the
same direction for a combination of the factors. Most contigs DE
according to the main effects of host-species or worm-population were
also DE according to the sex of the worm. There was also a number of
contigs differing for all three predictors in the same way. No contigs
were observed DE in both the host-species and worm-populations in the
same direction but not according to worm-sex. From the 68 contigs DE
in different \textit{A. crassus}-populations, 38 were also DE
according to worm sex and 16 according to all 3 main effects (see
figure \ref{venn_mod_man}).

The benefit of also allowing contrasting significant differences in
interaction terms highlights the power of the GLM-approach. In these
interactions a difference according to both focal factors in different
directions for factor combinations is indicated. For interactions
between host-species and parasite-population (eel*pop) for example
this mirrors the result of adult recovery i.e. a differential
regulation according to host-species/parasite-population combinations
found in nature.

And indeed also interaction-effects were observed: 7 contigs (0 OC)
showed differential expression according to the worm-sex*eel-species
interaction, 12 (3 OC) to worm-sex*parasite-population, 13 (2 OC) to
host-species*parasite-population, 1 (0 OC) contig showed significance
for the 3-way interaciton (see figure \ref{venn_mod_man}). It should
be noted, that conclusions drawn from of simple main effects do not
necessarily hold for contigs with significant interaction effects
(e.g. significantly higher expression in European population can then
mean higher only in \textit{An. anguilla}).

\figuremacroW{venn_mod_man}{Venn diagramm of contigs significant for
  different terms in \texttt{edgeR}-GLMs}{Overlap between differences
  in simple main effects are given as black numbers in the
  Venn-Diagramm. Numbers outside the circles in the lower left corner
  idicated non-significant contigs. The number of significant contigs
  for interaction effects are indicated in red for comparison. In (a)
  values for all contigs are given in (b) for orhtolog-confirmed (OC)
  contigs.}{0.6}

In summary, a low amount of overlap in main effects between
populations and host-species compared to the other main-effect
overlaps and in relation a higher proportion of interaction effects
between these two conditions was observed.

\section{Confirmation of contig categories through multivariate
  clustering}

We performed constrained redundancy analysis for the effects of
eel-host and worm-population. This technique similarly to principal
components analysis can partition the variance into orthogonal
components, and additionally constrain one of the components to the
factor of interest. We found 7\% of the variance in contigs DE between
eel-hosts and 11\% of the variance in contigs DE between
worm-population were explained by the corresponding factor. In both
evaluations more than 50\% of the remaining variance could be
explained by a single principal component, which was mainly consisting
of worm-sex (see figure \ref{pca_eel} a and \ref{pca_pop} a). When
only OC-DE contigs were considered the explained variance for
difference between eel-host droped to 3.3\% and the explained variance
for differences between worm-population was raised to 23\%, while the
sex-effect explained 70\% and 50\% of the variance (see figure
\ref{pca_eel} b and \ref{pca_pop} b). Significance of the constrained
component reported ba a premutation could be established at a p = 0.05
threshold for all but the OC eel-host DE subset.

\figuremacroW{pca_eel}{Constrained redundancy analysis for host-DE
  contigs}{Eel-host differences are displayed as constrained component
  on the x-axis, the principal component on the y-axis corresponds to
  the sex of the worm. (a) Host differences patition the variance in
  samples in like expcted for all contigs, the constrained component
  showed significance. (b) For OC contigs the constrained component
  fails to to partition the variance as expected, the component showed
  no significance for this subset of the data.}{0.60}

\figuremacroW{pca_pop}{Constrained redundancy analysis for
  population-DE contigs}{Population differences are displayed as
  constrained component on the x-axis, the principal component on the
  y-axis corresponds to the sex of the worm. Host differences patition
  the variance in samples like expcted for all contigs (a) as well as
  for OC-contigs (b). The constrained component showed significance in
  both subsets.}{0.60}

\afterpage{\clearpage}

\section{Biological processes associated with DE contigs}

We employed tests for overrepresentation of categories in
gene-ontology (GO). These tests respect the structure of the ontology
and also consider overrepresentation of higher level (ancestor-)
terms. Summarizing annotations at higher levels it is therefore
possible to conceive higher-order responses to the conditions
investigated.

For the differences between male and female worms enriched annotations
can be summarized into three broad categories: Terms overrepresented
due to spermatogenesis (eg. PP1-phosphatase and ester hydrolase are
important for speratogenesis in \textit{C. elegans}
\cite{wormbook_sperm, fardilha2011protein}) embryo development (many
obviouse terms) and terms for other processes putatively more related
other metabloic differences between males and females (such as
oxidoreductase activity; see table \ref{func-sex} but also additional
figures \ref{tGO_sex_BP_classic_10_all},
\ref{tGO_sex_CC_classic_10_all} and \ref{tGO_sex_MF_classic_10_all}).

\begin{longtable}{lp{4cm}rrrl}
\hline
GO.ID & Term & Annotated & Significant & Expected & p-value \\ 
\endfirsthead
\multicolumn{6}{c}%
{{\bfseries \tablename\ \thetable{} -- continued from previous page}} \\
\hline
GO.ID & Term & Annotated & Significant & Expected & p-value \\ 
\hline 
\endhead
\hline
\multicolumn{6}{|r|}{{Continued on next page}} \\ 
\hline
\endfoot
\endlastfoot
\hline
  \multicolumn{6}{l}{Molecular function} \\ 
  GO:0042578 & phosphoric ester hydrolase activity &  99 &  59 & 31.99 & 1.2e-08 \\ 
  GO:0016791 & phosphatase activity &  88 &  53 & 28.44 & 4.2e-08 \\ 
  GO:0004721 & phosphoprotein phosphatase activity &  65 &  42 & 21.00 & 6.5e-08 \\ 
  GO:0004722 & protein serine/threonine phosphatase act... &  34 &  24 & 10.99 & 4.8e-06 \\ 
  GO:0005509 & calcium ion binding &  78 &  43 & 25.21 & 2.1e-05 \\ 
  GO:0046873 & metal ion transmembrane transporter acti... &  32 &  21 & 10.34 & 0.00010 \\ 
  GO:0003824 & catalytic activity & 1354 & 482 & 437.55 & 0.00015 \\ 
  GO:0016614 & oxidoreductase activity, acting on CH-OH... &  46 &  27 & 14.86 & 0.00018 \\ 
  GO:0016616 & oxidoreductase activity, acting on the C... &  42 &  25 & 13.57 & 0.00023 \\ 
  GO:0017018 & myosin phosphatase activity &  10 &   9 & 3.23 & 0.00027 \\ 
  \hline
  \multicolumn{6}{l}{Biological process}  \\ 
  GO:0050896 & response to stimulus & 1535 & 583 & 504.78 & 1.7e-10 \\ 
  GO:0006470 & protein dephosphorylation &  63 &  41 & 20.72 & 1.2e-07 \\ 
  GO:0007391 & dorsal closure &  32 &  25 & 10.52 & 1.7e-07 \\ 
  GO:0016476 & regulation of embryonic cell shape &  13 &  13 & 4.27 & 5.0e-07 \\ 
  GO:0001700 & embryonic development via the syncytial ... &  49 &  33 & 16.11 & 6.7e-07 \\ 
  GO:0007392 & initiation of dorsal closure &  15 &  14 & 4.93 & 1.7e-06 \\ 
  GO:0046664 & dorsal closure, amnioserosa morphology c... &  15 &  14 & 4.93 & 1.7e-06 \\ 
  GO:0016311 & dephosphorylation &  86 &  49 & 28.28 & 2.6e-06 \\ 
  GO:0042221 & response to chemical stimulus & 864 & 337 & 284.12 & 3.1e-06 \\ 
  GO:0007394 & dorsal closure, elongation of leading ed... &  11 &  11 & 3.62 & 4.7e-06 \\ 
  \hline
  \multicolumn{6}{l}{Cellular compartment}  \\ 
  GO:0031224 & intrinsic to membrane & 372 & 164 & 118.85 & 8.4e-08 \\ 
  GO:0016021 & integral to membrane & 368 & 162 & 117.58 & 1.2e-07 \\ 
  GO:0005576 & extracellular region & 250 & 115 & 79.88 & 7.7e-07 \\ 
  GO:0031226 & intrinsic to plasma membrane & 176 &  86 & 56.23 & 1.0e-06 \\ 
  GO:0005887 & integral to plasma membrane & 172 &  84 & 54.95 & 1.4e-06 \\ 
  GO:0030054 & cell junction & 145 &  72 & 46.33 & 3.9e-06 \\ 
  GO:0000267 & cell fraction & 435 & 179 & 138.98 & 6.4e-06 \\ 
  GO:0016020 & membrane & 1154 & 417 & 368.70 & 3.6e-05 \\ 
  GO:0000164 & protein phosphatase type 1 complex &  14 &  12 & 4.47 & 4.9e-05 \\ 
  GO:0072357 & PTW/PP1 phosphatase complex &  14 &  12 & 4.47 & 4.9e-05 \\ 
  \hline\\
\caption[GO-terms enriched in DE between male and
female]{\textbf{GO-terms enriched in DE between male and female worms
    -} The top 10 enriched GO-categories are given for genes DE
  between the different male and female worms.}
\label{func-sex}
\end{longtable}

For the lower number of contigs DE between host-species obviousely
inferrence of higher order terms was only limited possible and in part
also unnecessary, because annotations can be interpreted at
face. However, annotations for contigs DE between eel-hosts
highlighted redundant terms associated with ``antigen processing and
presentation'' proteins which are in mammals usually involved in
antigen processing and cleavage of the invariant chain of the MHCII
complex. These terms led to Contig566 and Contig26 and thier
\textit{B. malayi}-orthologs ``Aspartic protease BmAsp-1, identical''
and ``Eukaryotic aspartyl protease family protein''. In blood feeding
helminths these enzymes are rather involved in early cleavage events
during the digestion of host haemoglobin \cite{pmid12782060}.

For contigs DE between worm populations despite the small 

\figuremacro{tGO_pop_BP_classic_10_all}{GO biological process graph
  for enriched terms in DE according to worm-population}{Subgraph of
  the GO-ontology cellular comparment category induced by the top 10
  terms identified as enriched in DE genes between different parasite
  populations. Boxes indicate the 10 most significant terms. Box color
  represents the relative significance, ranging from dark red (most
  significant) to light yellow (least significant). In each node the
  categoy-identifier, a (eventualy truncated) description of the term,
  the significance for enrichment and the number of DE / total number
  of annotated gene is given. Black arrows indicate a is ``is-a''
  relationship.}


\figuremacro{pop_all_heat}{Variance/mean stabilized expression values
  for contigs different between populations}{green indicates
  expression below the mean, red above the mean}




\section{Single gene differences}
\label{sec:single-gene-diff}










%%% Local Variables: ***
%%% mode:latex ***
%%% TeX-master: "../thesis.tex"  ***
%%% tex-main-file: "../thesis.tex" ***
%%% End: ***


% this file is called up by thesis.tex
% content in this file will be fed into the main document

\chapter{Discussion} % top level followed by section, subsection


% ----------------------- paths to graphics ------------------------

% change according to folder and file names
\ifpdf
    \graphicspath{{7/figures/PNG/}{7/figures/PDF/}{7/figures/}}
\else
    \graphicspath{{7/figures/EPS/}{7/figures/}}
\fi


% ----------------------- contents from here ------------------------

\section{Sanger-method pilot-sequencing}

One reason to sequence the libraries from the eels host was to
elucidate whether this contamination was nematode or species-typical
(e.g caused by poly-dT primers binding to A-rich rRNA regions), or
caused by shortcomings in the preparation.

Nevertheless the stringent quality trimming and processing of raw
reads, as summarized in the present document, make the remaining ESTs
a valuable resource for comparison with future 454-sequencing-data..




\section{454-pyrosequencing}




%%% Local Variables: ***
%%% mode:latex ***
%%% TeX-master: "../thesis.tex"  ***
%%% tex-main-file: "../thesis.tex" ***
%%% End: ***
               % discussion of results


% this file is called up by thesis.tex
% content in this file will be fed into the main document

\chapter{Materials \& methods} % top level followed by section, subsection


% ----------------------- paths to graphics ------------------------

% change according to folder and file names
\ifpdf
    \graphicspath{{8/figures/PNG/}{8/figures/PDF/}{8/figures/}}
\else
    \graphicspath{{8/figures/EPS/}{8/figures/}}
\fi

% ----------------------- contents from here ------------------------


\section{Sampling of worms from wild eels for Sanger- and
  pyrosequencing}

Cultured eels were acquired from an aquaculture directly adjacent to
Kaoping river (22.6418N; 120.4440E) 15km stream upwards from its
estuary. Wild eel were bought from a fisherman, fishing in the estuary
of Kao-Ping river (22.5074N; 120.4220E). All eels were transported to
the Institute of Fisheries Science at the National Taiwan University
in Taipei in aerated plastic bags, where they were sheltered until
dissection.

Eels were decapitated, length (to the nearest 1.0mm) and weight (to
the nearest 0.1g) were measured, and sex was determined by visual
inspection of the gonads. The swimbladder was opened, adult worms were
removed from the lumen with a forceps, their sex was determined, and
they were counted. All adult \textit{A. crassus} were preserved in
RNAlater(Quiagen, Hilden, Germany) in individual plastic tubes.

Worms from the European eel were sampled in Sniardwy Lake, Poland
(53.751959N; 21.730957E) by Urszula Weclawski and from the
Linkenheimer Altrhein, Germany (49.0262N; 8.310556E), following a
procedure similar to the one described above for worms from Taiwan.

\section{RNA-extraction and cDNA synthesis for Sanger- and
  pyrosequencing}

Total RNA was extracted from single, whole worms using the RNeasy kit
(Quiagen, Hilden, Germany), following the manufacturers
protocol. Alternatively parts of the liver of the host species
\textit{Anguilla japonica}, which also had been preserved in RNAlater
were used for RNA extraction, following the same protocol.

The Evrogen MINT cDNA synthesis kit (Evrogen, Moscow, Russia ) was
then used to amplify mRNA transcripts according to the manufacturers
protocol. It uses an adapter sequence at 3' the end of a poly
dT-primer for first strand synthesis and adds a second adapter
complementary to the bases at the 5' end of the transcripts by
terminal transferase activity and template switching. Using these
adapters it is possible to specifically amplify mRNA enriched for
full-length transcripts.

\section{Cloning for Sanger-sequencing}

The obtained cDNA preparations were undirectionally cloned into
TOPO2PCR-vectors (Invitrogen, Carlsbad, USA) and TOP10 chemically
competent cells (Invitrogen, Carlsbad, USA) were transformed with this
construct. The cells were plated on LB-medium-agarose containing
Kanamycin (5mg/ml), xGal
(5-bromo-4-chloro-3-indolyl-$\beta$-D-galactopyranoside) and IPTG
(Isopropyl-$\beta$-D-1-thiogalactopyranosid). After 24h of incubation
at $36\,^{\circ}\mathrm{C} $ cells were picked into 96-well
micro-liter-plates containing liquid LB-medium and Kanamycin (5mg/ml)
and incubated for another 24h. Subsequently 2ml of the cells were used
as template for amplification of the insert by PCR using the primers
\begin{description}
\item[Forward] M13F(GTAAAACGACGGCCAGT) and
\item[Reverse] M13R(GGCAGGAAACAGCTATGACC)
\end{description}
in a concentration of 10$\mu$M. The protocol for PCR cycling is shown
% in table \ref{tab:PCR}.

\begin{table}[h]
  \centering
  \begin{tabular}{lllll} 
    \textbf{Initial denaturation} &  $ 94\, ^{\circ}\mathrm{C} $ & 5min &  &\\ 
    \hline
    \textbf{Denaturation} &  $ 94\, ^{\circ}\mathrm{C} $ &30s& & \\ 
    \textbf{Annealing} &   $ 54\, ^{\circ}\mathrm{C} $ & 45s & 35 cycles &\\ 
    \textbf{Elongation} &   $ 72\, ^{\circ}\mathrm{C} $ & 2min &  &\\ 
    \hline
    \textbf{Final Elongation} &   $ 72\, ^{\circ}\mathrm{C} $ & 10min &\\ 
  \end{tabular}   
  \caption{PCR protocol for insert amplification}
  \label{tab:PCR}
\end{table}

Amplification products were controlled on gel and cleaned using SAP
(Shrimp Alkaline Phosphatase) and ExoI (Exonuclease I). Sequencing
reactions were performed using the BigDye-Terminator kit and
PCR-primers (forward or reverse) in a concentration of 3.5$\mu$M and
sequenced on an ABI 3730 DNA Analyzer (Applied Biosystems, Foster
City, California, USA).  For \textit{A. crassus} the following
libraries were prepared:
 
\begin{description}
\item{Ac\_197F:} Female from Taiwanese aquaculture
\item{Ac\_106F:} Female from Taiwanese aquaculture
\item{Ac\_M175:} Male from Taiwanese aquaculture
\item{Ac\_FM:} Female from Taiwanese aquaculture
\item{Ac\_EH1:} Same cDNA preparation as Ac\_FM, but sequenced by
  students in a practical
\end{description}

For \textit{Anguilla japonica} the following three libraries:
\begin{description}
\item{Aj\_Li1:} liver of an eel from aquaculture
\item{Aj\_Li2:} liver of an eel from aquaculture
\item{Aj\_Li3:} liver of an eel from aquaculture
\end{description}

\section{Pilot Sanger-sequencing}

The original sequencing-chromatographs ("trace-files") were renamed
according to the NERC environmental genomics scheme. "Ac" was used as
project-identifier for \textit{Anguillicola crassus}, "Aj" for
\textit{An. japonica}. In \textit{Anguillicola} sequences information
on the sequencing primer (forward or reverse PCR primer;
\textit{An. japonica} sequences were all sequenced using the forward
PCR primer) was stored in the middle "library"-field, resulting in
names of the following form:

\begin{itemize}
\item \texttt{Ac\_[\textbackslash{}d|\textbackslash{}w]\{2,4\}(f|r)\_\textbackslash{}d\textbackslash{}d\textbackslash{}w\textbackslash{}d\textbackslash{}d}
\item  \texttt{Aj\_[\textbackslash{}d|\textbackslash{}w]\{2,4\}\_\textbackslash{}d\textbackslash{}d\textbackslash{}w\textbackslash{}d\textbackslash{}d}
\end{itemize}

The last field indicates the plate number (two digits), the row (one
letter) and the column (two digits) of the corresponding clone. For
first quality trimming \texttt{trace2seq}, a tool derived from
\texttt{trace2dbEST} (both part of \texttt{PartiGene}
\cite{parkinson_partigene--constructing_2004}) was used, briefly it
performs quality trimming using \texttt{phred}
\cite{ewing_base-calling_1998} and trimming of vector sequences using
\texttt{cross-match} \cite{PHRAP}. The adapters used by the MINT kit
were trimmed by supplying them in the vector-file used for trimming
along with the TOPO2PCR-vector. After processing with
\texttt{trace2seq} additional quality trimming was performed on the
produced sequence-files using a custom script. This trimming was
intended to remove artificial sequences produced when the sequencing
reaction starts at the 3' end of the transcript at the poly-A
tail. These sequences typically consist of numerous homo-polymer-runs
throughout their length caused by "slippage" of the reaction.
The basic Perl regular expression used for this was:\\

\texttt{/(.*A\{5,\}|T\{5,\}|G\{5,\}|C\{5,\}.*)\{\$lengthfac,\}/g}\\

Where \texttt{\$lengthfac} was set to the length of the sequence
divided by 70 and rounded to the next integer. So only one
homo-polymer-run of more then 5 bases was allowed per 75 bases.

Sequences were screened for host contamination by a comparison of
\texttt{Blast} searches against nempep \cite{pmid21550347} (version 4)
and a fish protein database. Sequences producing better bit scores
against fish proteins than nematode proteins were labeled as
host-contamination.

Only the trace-files corresponding to the sequences still regarded as
good after this step were processed with
\texttt{trace2dbEST}. Additionally to the processing of traces already
included in \texttt{trace2seq} sequences were preliminary annotated
using \texttt{Blast} versus the NCBI-NR non-redundant protein database
and EST-submission-files were produced.

\section{Pyrosequencing}

\subsection{cDNA preparation and sequencing}

RNA was extracted from individual adult male and female nematodes and
from a population of L2 larvae. RNA was reverse transcribed and
amplified into cDNA using the MINT-cDNA synthesis kit (Evrogen,
Moscow, Russia).  For host contamination screening a liver-sample from
an uninfected \textit{An. japonica} was also processed. Emulsion PCR
was performed for each cDNA library according to the manufacturer's
protocols (Roche/454 Life Sciences), and sequenced on a Roche 454
Genome Sequencer FLX. All samples were sequenced using the FLX
Titanium chemistry, except for the Taiwanese female sample T2, which
was sequenced using FLX standard chemistry, to generate between 99,000
and 209,000 raw reads. For the L2 larval library, which had a larger
number of non-\textit{A. crassus}, non-\textit{Anguilla} reads,
screening Roche 454 data produced on the same run in independent
sequencing lanes confirmed that these data were not laboratory
contaminants.

\subsection{Trimming, quality control and assembly}

Raw sequences were extracted in \texttt{fasta}-format (with the
corresponding qualities files) using \texttt{sffinfo} (Roche/454) and
screened for adapter sequences of the MINT-amplification-kit using
\texttt{cross-match} \cite{PHRAP} (with parameters \texttt{-minscore
  20 -minmatch 10}). \texttt{Seqclean} \cite{tgicl_pertea} was used to
identify and remove poly-A-tails, low quality, repetitive and short
(<100 base) sequences. All reads were compared to a set of screening
databases using \texttt{Blast} (expect value cutoff E<1e-5, low
complexity filtering turned off: -F F). The databases used were (a) a
host sequence database comprising an assembly of the
\textit{An. japonica} Roche 454 data, an unpublished assembly of
\textit{An. anguilla} Sanger dideoxy sequenced expressed sequence tags
(made available to us by Gordon Cramb, University of St Andrews) and
transcripts from EeelBase \cite{pmid21080939} a publicly available
transcriptome database for the European eel; (b) a database of
ribosomal RNA (rRNA) sequences from eel species derived from my Roche
454 data and EMBL-Bank; and (c) a database of rRNA sequences
identified in my \textit{A. crassus} data by comparing the reads to
known nematode rRNAs from EMBL-Bank. This last database notably also
contained xenobiont rRNA sequences. Reads with matches to one of these
databases over more than 80\% of their length and with greater than
95\% identity were removed from the dataset. Screening and trimming
information was written back into sff-format using \texttt{sfffile}
(Roche 454). The filtered and trimmed data were assembled using the
combined assembly approach \cite{pmid20950480}: two assemblies were
generated, one using \texttt{Newbler v2.6} \cite{pmid16056220} (with
parameters \texttt{-cdna -urt}), the other using \texttt{Mira v3.2.1}
\cite{miraEST} (with parameters
\texttt{--job=denovo,est,accurate,454}). The resulting two assemblies
were combined into one using \texttt{Cap3} \cite{Cap3_Huang} at
default settings and contigs were labeled by whether they derived from
both assemblies or one assembly only.

\subsection{Evaluation of the assemblies}

The ace-files for all three (two first-order, one second-order)
assemblies were interrogated for the fate of single reads. This was
used to tabulate the full read-first-order-second-order-associations.

\texttt{Blast} (\texttt{blastx -e 1e-5}) was used to search the
complete proteomes of \textit{C. elegans} (as present in wormbase
v.220) and the complete proteome of \textit{B. malayi} (as present in
uniref 100) for the contigs and singletons of all investigated
assemblies. A custom Perl-script (provided by S. Kumar) was then used
to mask all bases in the database covered. For each sequence in the
database the size of the masked region was then determined and
statistics were created summarising the number of database-sequences
with any coverage, the number with coverage over 80\% of their
sequence-length and the overall proportion of bases covered.

Based on reads shared between clusters I collapsed reads linked by
such read-paths, assigned a cluster-id and recorded the size of the
cluster.

To estimate contig-coverage I converted \texttt{sam}-output generated
with \texttt{ssaha2}\cite{pmid11591649} via a sorted \texttt{bam}-file
to \texttt{pileup}-format using \texttt{samtools}
\cite{journals/bioinformatics/LiHWFRHMAD09}. For a second evaluation
I excluded best-hits mapping to multiple contigs before converting the
\texttt{sam}-file to obtain unique coverage.

\subsection{Post-assembly classification and taxonomic assignment of
  contigs}

After assembly contigs were assessed a second time for host and other
contamination by comparing them (using \texttt{Blast}) to the three
databases defined above, and also to nembase4, a nematode
transcriptome database derived from whole genome sequencing and EST
assemblies \cite{parkinson_nembase:resource_2004, pmid21550347}. For
each contig, the highest-scoring match was recorded as long as it
spanned more than 50\% of the contig. I also compared the contigs to
the NCBI non-redundant nucleotide (NCBI-nt) and protein (NCBI-nr)
databases, recording the taxonomy of all best matches with expect
values better than 1e-05. TUGs with a best hit to non-Metazoans and to
Chordata within Metazoa were additionally excluded from further
analysis.

\subsection{Protein prediction and annotation}
\label{pro4est}

Protein translations were predicted from the contigs using
\texttt{prot4EST} (version 3.0b)
\cite{wasmuth_prot4est:_2004}. Proteins were predicted either by
joining single high scoring segment pairs (HSPs) from a \texttt{Blast}
search of uniref100 \cite{pmid18836194}, or by \texttt{ESTscan}
\cite{estscan}, using as training data the \textit{Brugia malayi}
complete proteome back-translated using a codon usage table derived
from the \texttt{Blast} HSPs, or, if the first two methods failed,
simply the longest ORF in the contig. For contigs where the protein
prediction required insertion or deletion of bases in the original
sequence, I also imputed an edited sequence for each affected
contig. Annotations with Gene Ontology (GO), Enzyme Commission (EC)
and Kyoto Encyclopedia of Genes and Genomes (KEGG) terms were inferred
for these proteins using \texttt{Annot8r} (version 1.1.1)
\cite{schmid_annot8r:_2008}, using the annotated sequences available
in uniref100 \cite{pmid18836194}. Up to 10 annotations based on a
\texttt{Blast} similarity bitscore cut-off of 55 were obtained for
each annotation set. The complete \textit{B. malayi} proteome (as
present in uniref100) and the complete \textit{C. elegans} proteome
(as present in wormbase v.220) were also annotated in the same
way. \texttt{SignalP V4.0} \cite{pmid21959131} was used to predict
signal peptide cleavage sites and signal anchor signatures for the
\textit{A. crassus}-transcriptome and similarly again for the
proteomes of the tow model-worms.  Additionally \texttt{InterProScan}
\cite{pmid11590104} (command line utility \texttt{iprscan} (version
4.6) with options \texttt{-cli -format raw -iprlookup -seqtype p
  -goterms}) was used to obtain domain based annotations for the high
credibility assembly (highCA) derived contigs.

I recorded the presence of a lethal RNAi-phenotype in the
\textit{C. elegans} ortholog of each TUG using the biomart-interface
\cite{pmid22083790} to wormbase v. 220 through the R-package
\texttt{biomaRt} \cite{pmid19617889}.

\subsection{Single nucleotide polymorphism analysis}

I mapped the raw reads against the the complete set of contigs,
replacing imputed sequences for originals where relevant, using
\texttt{ssaha2} \cite{pmid11591649} (with parameters \texttt{-kmer 13
  -skip 3 -seeds 6 -score 100 -cmatch 10 -ckmer 6 -output sam -best
  1}). From the \texttt{ssaha2} output, pileup-files were produced
using \texttt{samtools} \cite{journals/bioinformatics/LiHWFRHMAD09},
discarding reads mapping to multiple regions. \texttt{VarScan}
\cite{pmid19542151} (\texttt{pileup2snp}) was used with default
parameters on pileup-files to output lists of single nucleotide
polymorphisms (SNPs) and their locations. For enrichment analysis of
GO-terms I used the R-package \texttt{GOstats} \cite{pmid17098774}.

Using \texttt{Samtools} \cite{journals/bioinformatics/LiHWFRHMAD09}
(\texttt{mpileup -u}) and \texttt{Vcftools} \cite{pmid21653522}
(\texttt{view -gcv}) I genotyped individual libraries for the list of
previously found overall SNPs. Genotype-calls were accepted at a
phred-scaled genotype quality threshold of 10. In addition to the
relative heterozygosity (number of homozygous sites/number of
heterozygous sites) I used the R package \texttt{Rhh}
\cite{pmid21565077} to calculate internal relatedness
\cite{pmid11571049}, homozygosity by loci \cite{pmid17107491} and
standardised heterozygosity \cite{coltman81j} from these data.

Using 1000 bootstrap replicates I confirmed the significance of
heterozygote-heterozygote correlation by analysing the mean and 95\%
confidence intervals from 1000 bootstrap replicates estimated for all
measurements.

\subsection{Gene-expression analysis}

Read-counts were obtained from the \texttt{bam}-files generated also
for genotyping using the R-package \texttt{Rsamtoools}
\cite{rsamtools}. TUGs with less than 48 reads over all libraries were
excluded from analysis, as diagnostic plot (not shown) indicated a
lack of statistical power for lower overall expression. I used the
R-package \texttt{DESeq} \cite{pmid20979621} (version 1.6.1) to assess
statistical significance of differences in counts according to groups
of libraries.

Additionally I collapsed TUGs by their orthologous assignment in
\textit{C. elegans} and \textit{B. malayi}. I used the sums of counts
for these orthologous-groups to asses the influence of mapping to my
potentially fragmented reference. For both model-nematodes fold-change
and p-values were obtained the same way than for the contigs and
merged with these.

\subsection{Enrichment analysis}
\label{GO-over-meth}

We used Mann-Whitney u-tests to test the influence of factors on dn/ds
values, when multiple contrasts between groups (factors) were
investigated we used Nemenyi-Damico-Wolfe-Dunn tests. For
overrepresentation of one group (factor) in another group (factor) we
used Fisher's exact test.

Prior to analysis of GO-term over-representation (based on dn/ds or
expression values) we used the R-package \texttt{annotationDbi}
\cite{AnnotationDbi} to obtain a full list of associations (also with
higher-level terms) from \texttt{annot8r}-annotations. We then used the
R-package \texttt{topGO} \cite{topGO} to traverse the annotation-graph
and analyse each node in the annotation for over-representation of the
associated term in the focal gene-set compared to an appropriate
universal gene-set (all contigs with dn/ds values or all contigs
analysed for gene-expression) with the ``classic'' method and Fisher's
exact test.

\section{Transcriptomic divergence in a common garden experiment}

\subsection{Experimental infection of eels}

\textit{An. anguilla} were obtained from the Albe-Fishfarm in
Haren-R\"utenbrock, Germany. \textit{An. japonica} were caught at the
glass-eel stage in the estuary of Kao-ping River, Taiwan by
professional fishermen and kept at a water temperature of
26$^{\circ}$C until they reached a size of $>$ 35 cm.

The absence of infections with \textit{A. crassus} in both eel-species
was confirmed by dissection of 10 individuals of each species.

After an acclimatisation period of 4 weeks (\textit{An. anguilla}) or
when they reached a size of $>$ 35cm (\textit{An. japonica}) eels were
infected using a stomach tube as described in
\cite{boon1990effect}. During the infection period water temperature
was held constant at 20$^{\circ}$C. Eels were kept in 160-litre tanks
in groups of 5-10 individuals and continuously provided with fresh,
oxygenated water and once every two days with commercial fish pellets
(Dan-Ex 2848, Dana Feed A/S Ltd, Horsens, Denmark) \textit{at
  libitum}.

L2 larvae used for the infection were collected from the swimbladders
of wild yellow and silver eels from the River Rhine near Karlsruhe and
from Lake M\"uggelsee near Berlin in Germany. Taiwanese larvae were
obtained from eels from an aquaculture adjacent to Kao Ping River in
south Taiwan and from a second aquaculture in Yunlin county,
approximately 150 km further north on the west coast of Taiwan. They
were stored at 4$^{\circ}$C for no longer than 2 weeks before copepods
were infected. Mixed species samples of uninfected copepods were
collected from a small pond near Karlsruhe, known to be free of eels
(and \textit{Anguillicola}). They were infected individually in wells
of micro-titer plates at an intensity of roughly 10 L2-larvae per
copepod. One week after infection they were placed in bigger
tanks. Twice a week yeast was provided as food and at 21 dpi infective
L3 were harvested with using a tissue potter as described by
\cite{haenen_improved_1994}. 50 L3 for infection of individual eel
were suspended in 100 $\mu$l RPMI-1640 medium (Quiagen, Hilden,
Germany) and eels were infected as described above.

55-57 days post infection (dpi) eels were euthanized and dissected.
The swimbladder was opened and after determination of the sex of adult
worms under a binocular microscope (Semi 2000, Zeiss, Germany), adult
\textit{A. crassus} were immediately immersed in RNAlater (Quiagen,
Hilden, Germany).

\subsection{RNA extraction and preparation of sequencing libraries}

RNA was extracted from 12 individual female worms and for 12 pools of
male worms using the RNeasy-kit (Quiagen, Hilden, Germany) (see table
\ref{tab:lib-prep}).

The paired-end TruSeqTM RNA sample preparation kit (Illumina) was
followed to build sequencing libraries with insert sizes of roughly
270 bp for paired-end sequencing from cDNA libraries: briefly, poly-T
oligo-attached magnetic beads were used for purification of mRNA and
to simultaneously fragment the RNA. The RNA was then primed with
random hexamer primers for cDNA synthesis and reverse transcribed into
first strand cDNA using reverse transcriptase. The cDNA was cleaned
from the second strand reaction, overhangs were repaired to form blunt
ends, a single ``A''-nucleotide was added at the 3' end and paired end
sequencing adapters ware ligated with a complementary
``T''-overhang. In this step multiple differently indexed paired-end
adapters were used to enable multiplexing of the 24 different
sequencing libraries in 3 pools of 8 samples each. These three pools
all contained one random replicate each for each treatment combination
ensuring complete statistical independence of replicates. Molecules
having adapter sequences were enriched in the mix using PCR and the
libraries were controled for quality and quantity on the BioAnalyzer
(Agilent). Clusters were generated by bridge amplification. The
resulting clusters were sequenced on the Genome Analyzer IIX in
combination with the paired-end module. The first read was sequenced
using using the first primer Rd1 SP. The original template strand was
then used to regenerate the complementary strand, the original strand
was removed and complementary strand acted as a template for the
second read, sequenced primed by the second sequencing primer Rd2 SP.

\subsection{Mapping and normalisation of read-counts}

All sequencing reads were mapped to the fullest 454 assembly (as
defined in \ref{sec:final-full-assembly}; I were including TUGs
inferred as host or xenobiont origin as filter) using \texttt{BWA}
\cite{pmid20080505} (version 0.5.9-r16; \texttt{BWA aln} and
\texttt{BWA sampe} with default options) and processed with
\texttt{samtools} \cite{journals/bioinformatics/LiHWFRHMAD09} (version
0.1.18; \texttt{samtools view -uS -q 1}) to only allow uniquely
mapping reads. All reads mapping to host- and xenobiont off-target
data were removed during downstream evaluation.

Counts were summed for technical replicates and counts to
lowCA-derived contigs were disregarded for statistics on a
contigs-base as well as spurious read counts to contigs with less
than 32 mapping reads in total (see however \ref{collapse} for how
these counts were used in further tests of reference fragmentation).

The remaining counts were normalised using \texttt{DESeq} (version
1.6.1) (i.e.  the normalisation factor was estimated by the median of
scaled counts, similar to the weighted trimmed mean of the log
expression ratios used later in \texttt{edgeR}). All tables
summarising read-counts are based on these normalised counts. I
obtained ``variance stabilised data'' in an expression matrix for each
gene and library using the ``blind'' option in a calculation not
informed (and biased by) the model-design. These data were used in all
gene-centring heatmap and multivariate visualisations. Additionally this
matrix was transposed to get sample-to-sample distances.

\subsection{Statistical analysis with generalised linear models (GLMs)}

The R-package \texttt{edgeR} (version 2.4.1) \cite{pmid19910308} was
used to build negative binomial generalised linear models, as these
specialised GLMs outperformed GLMs in \texttt{DESeq} in speed and
reliability of convergence. Modeled were based on a negative binomial
distribution and the dispersion parameter for each transcript was
approximated with a trend depending on the overall level of
expression. In the maximal fitted model expression was regressed on
worm-sex, host-species and parasite population, including all their
interactions. The full model thus contained terms $ S_i + H_j + P_k +
(SH)_{ij} + (SP)_{ik} + (HP)_{jk} + (SHP)_{ijk} + \varepsilon$, where
$\varepsilon$ is the residual variance, $S_i$ is the effect of the ith
sex (male or female), $H_j$ is the effect of the ith host species
(\textit{An. anguilla} or \textit{An. japonica}), $P_k$ is the effect
of the kth population (European or Asian), $(SH)_{ij}$ is the
sex-by-species interaction and similarly for the other interactions.

The hierarchical nature of generalised linear models was respected
considering (removing) all interaction effects of a main-term
(e.g. $(SP)_{ik}$, $(SH)_{ij}$ and $(SHP)_{ijk}$) when analysing
models for the significance of that term (e.g. $S_i$). Resulting
p-values were corrected for multiple testing using the method of
Benjamini and Hochberg \cite{benjamini1995controlling} and
differential expression was inferred at a false discovery rate (FDR)
of 5\% (adjusted p-value of 0.05).

\subsection{Count-collapsing for orthologs from two model-species}
\label{collapse}

In order to test the influence of deficiencies (i.e. fragmentation) of
the assembly on mapping and read-counts I summed read counts over
orthologous sequence in \textit{C. elegans} and
\textit{B. malayi}. For this purpose I used all reference contigs
(also lowCA-derived contigs to allow inference of fragmentary mapping
to those, but not contigs of non-\textit{A. crassus} origin).
Differential expression for these orthologous-counts was analysed the
same way as for contigs. Contigs were filtered based on inference
from orthologous counts merging the two orthologous evaluations and
the contig evaluation. Differential expression was accepted at a FDR
of 5\% for the contig evaluation and 10\% for both of the two
orthologous evaluations.

\subsection{Multivariate confirmation of linear models}
\label{sec:mult-conf-line}

I used the R-package \texttt{vegan} (version 2.0-2) to perform
constrained redundancy analysis on contigs identified as significant
in GLMs before. For each set of contigs (different for sex, eel-host
or worm-population) the appropriate constrained component was
used. The proportion of the variance explained by the constrained
component was recorded and the constrained component was tested for
significance using a permutation test implemented in \texttt{vegan}.

\subsection{GO-term enrichment analysis}

Enrichment analysis was performed as described above for
pyrosequencing data (see \ref{GO-over-meth}).

\subsection{Clustering analysis}

The R-package \texttt{HeatmapPlus} was used on variance stabilised
expression values to visualise hierarchical clusters similar to the
method of \cite{pmid9843981}. The results were displayed along with
annotations stored in a Bioconductor eSet-class object.

\section{General coding methods}

The bulk of analysis (unless otherwise cited) presented in this paper
was carried out in \texttt{R} \cite{R_project} using custom scripts.
I used a method provided in the R-packages \texttt{Sweave}
\cite{lmucs-papers:Leisch:2002} and \texttt{Weaver} \cite{weaver} for
``reproducible research'' combining \texttt{R} and \LaTeX code in a
single file. The complete reproducible compilations were only carried
out for sub-chapters of this document, the thesis-document was then
compiled from plain \LaTeX sub-documents. Nevertheless all
intermediate data files needed to compile sub-document of the thesis
from data-sources are provided upon request. For general visualisation
I used the R-packages \texttt{ggplot2} \cite{ggplot-book} and
\texttt{VennDiagram} \cite{pmid21269502}.


%%% Local Variables: ***
%%% mode:latex ***
%%% TeX-master: "../thesis.tex"  ***
%%% tex-main-file: "../thesis.tex" ***
%%% End: ***
 






        % description of lab methods



% --------------------------------------------------------------
%:                  BACK MATTER: appendices, refs,..
% --------------------------------------------------------------

% the back matter: appendix and references close the thesis


%: ----------------------- bibliography ------------------------

% The section below defines how references are listed and formatted
% The default below is 2 columns, small font, complete author names.
% Entries are also linked back to the page number in the text and to
% external URL if provided in the BibTex file.

% PhDbiblio-url2 = names small caps, title bold & hyperlinked, link to page 
\begin{multicols}{2} % \begin{multicols}{ # columns}[ header text][ space]
\begin{tiny} % tiny(5) < scriptsize(7) < footnotesize(8) < small (9)

\bibliographystyle{Latex/Classes/PhDbiblio-url2} % Title is link if provided
\renewcommand{\bibname}{References} % changes the header; default: Bibliography

\bibliography{/home/ele/bibtex/master,/home/ele/bibtex/master2} % adjust this to fit your BibTex file

\end{tiny}
\end{multicols}

% --------------------------------------------------------------
% Various bibliography styles exit. Replace above style as desired.

% in-text refs: (1) (1; 2)
% ref list: alphabetical; author(s) in small caps; initials last name; page(s)
%\bibliographystyle{Latex/Classes/PhDbiblio-case} % title forced lower case
%\bibliographystyle{Latex/Classes/PhDbiblio-bold} % title as in bibtex but bold
%\bibliographystyle{Latex/Classes/PhDbiblio-url} % bold + www link if provided

%\bibliographystyle{Latex/Classes/jmb} % calls style file jmb.bst
% in-text refs: author (year) without brackets
% ref list: alphabetical; author(s) in normal font; last name, initials; page(s)

%\bibliographystyle{plainnat} % calls style file plainnat.bst
% in-text refs: author (year) without brackets
% (this works with package natbib)


% --------------------------------------------------------------

% according to Dresden med fac summary has to be at the end
%
% Thesis Abstract -----------------------------------------------------


%\begin{abstractslong}    %uncommenting this line, gives a different abstract heading
\begin{abstractslong}        %this creates the heading for the abstract page

  The ability to expand into new environments and niches, despite
  being highly adapted to a habitual environment, is a fascinating
  feat of organisms. 30 years ago \textit{Anguillicola crassus} was
  introduced from Asia, where it parasitises \textit{Angilla
    japonica}, to Europe and spread here in the new host species
  \textit{Anguilla anguilla}. Whether and how much phenotypic
  plasticity or rapid adaptation to differential selection are
  contributing to its success in invading new host-populations is a
  question of substantial evolutionary interest.

  Gene regulatory networks, as an important link between genotype
  and phenotype, are thought to play a central role both in the
  response to stress (e.g. from as yet unexperienced environmental
  stressors) and in local adaptation.

  In the present project, differential gene-expression in
  \textit{A. crassus} populations was assessed using next generation
  sequencing on the 454 and Illumina platforms and genetic components
  of differences were isolated in cross-inoculation experiments with
  both Asian and European host-species and parasite populations.

  Several proteases were shown to be under positive selection on the
  sequence level, highlighting this group of enzymes as possible
  targets of an immune-attack on \textit{A. crassus}. On the
  gene-expression level, the extent of heritable change was large in
  comparison to the effect of modification in different
  host-environments. Mitochondrially encoded subunits of the
  respiratory chain and other genes connected to aerobic respiration
  showed divergent expression patterns in European vs. Asian parasite
  populations; cuticle collagen genes showed ``adapted'' patterns of
  expression in present day sympatric host-parasites pairs.
  
  These results identified gene-expression phenotypes, confirming the
  divergence of European \textit{A. crassus} populations. Such
  phenotypes will be more accessible to popluation-genetic analysis
  invetigating selection than complex life history traits.

\end{abstractslong}
%\end{abstractlongs}


% ---------------------------------------------------------------------- 

%%% Local Variables: ***
%%% mode:latex ***
%%% TeX-master: "../thesis.tex"  ***
%%% tex-main-file: "../thesis.tex" ***
%%% End: ***
     

%: Declaration of originality

% Thesis statement of originality -------------------------------------

% Depending on the regulations of your faculty you may need a declaration like the one below. This specific one is from the medical faculty of the university of Dresden.

\begin{declaration}        %this creates the heading for the declaration page

  I herewith declare that I have produced this paper without the
  prohibited assistance of third parties and without making use of
  aids other than those specified; notions taken over directly or
  indirectly from other sources have been identified as such. This
  paper has not previously been presented in identical or similar form
  to any other German or foreign examination board.

  Chapter \ref{cha:pyro} was in similar form submitted for publication
  to BMC Genomics, in the course of manuscript preparation Mark
  Blaxter edited the text.
  
  Chapter \ref{cha:tra-diff} is in similar intended for publication in
  Plos biology, Mark Blaxter edited parts of the text.

  The thesis work was conducted from May 2008 to December 2011 under
  the supervision of Prof. Dr. Horst Taraschewski at the Karlsruhe
  Institute of Technology and Prof. Mark Blaxter at the University of
  Edinburgh.

\vspace{10mm}

KARLSRUHE, 


\end{declaration}


%%% Local Variables: ***
%%% mode: latex ***
%%% TeX-master: "../thesis.tex" ***
%%% tex-main-file: "../thesis.tex" ***
%%% End: ***

% ----------------------------------------------------------------------



\end{document}
